\section*{Розділ I. Загальні положення}
\addcontentsline{toc}{section}{Розділ I. Загальні положення}

\subsection*{1.1. Статус та правові засади діяльності}
\addcontentsline{toc}{subsection}{1.1. Статус та правові засади діяльності}
    1.1.1. Центральна виборча комісія студентів Київського авіаційного інституту (далі -- ЦВКс) є постійно діючим органом студентського самоврядування, відповідальним за організацію та проведення виборів до органів студентського самоврядування Університету.

    1.1.2. ЦВКс діє відповідно до Конституції України, Закону України ``Про вищу освіту'', Статуту Університету, Положення про органи студентського самоврядування Київського авіаційного інституту (далі -- Положення про ОСС) та цього Положення.

    1.1.3. Це Положення діє відповідно до та на виконання Положення про ОСС і визначає порядок створення, організації діяльності, повноваження та функції ЦВКс.

\subsection*{1.2. Незалежність та підзвітність}
\addcontentsline{toc}{subsection}{1.2. Незалежність та підзвітність}
    1.2.1. ЦВКс є незалежним органом у питаннях організації та проведення виборів до органів студентського самоврядування та не підлягає втручанню з боку адміністрації Університету чи інших органів студентського самоврядування при здійсненні своїх виборчих повноважень.

    1.2.2. ЦВКс звітує про свою діяльність перед Конференцією студентів Університету (КСУ) шляхом подання річного звіту та звітів за результатами проведених виборів.

    1.2.3. КСУ здійснює контроль за діяльністю ЦВКс через заслуховування звітів, але не втручається у виборчі процедури, що проводяться ЦВКс.

\subsection*{1.3. Принципи діяльності}
\addcontentsline{toc}{subsection}{1.3. Принципи діяльності}
    1.3.1. Діяльність ЦВКс здійснюється на основі принципів:

        \begin{enumerate}[label=\alph*)]
            \item Законності -- строгого дотримання вимог законодавства України, Статуту Університету та нормативних актів органів студентського самоврядування;
            \item Незалежності -- самостійного прийняття рішень у межах своїх повноважень без втручання інших органів;
            \item Неупередженості -- забезпечення рівних можливостей для всіх кандидатів та виборців;
            \item Прозорості -- відкритості виборчого процесу для громадського контролю;
            \item Гласності -- забезпечення доступу до інформації про виборчий процес;
            \item Колегіальності -- прийняття рішень більшістю голосів членів ЦВКс;
            \item Відповідальності -- персональної та колегіальної відповідальності за прийняті рішення та проведені заходи.
        \end{enumerate}

\subsection*{1.4. Основні завдання та функції}
\addcontentsline{toc}{subsection}{1.4. Основні завдання та функції}
    1.4.1. Основними завданнями ЦВКс є:

        \begin{enumerate}[label=\alph*)]
            \item Забезпечення реалізації виборчих прав студентів Університету;
            \item Організація та проведення виборів до органів студентського самоврядування на засадах прямого таємного голосування відповідно до ст. 40 Закону України ``Про вищу освіту'';
            \item Забезпечення демократичності, законності та прозорості виборчого процесу;
            \item Сприяння підвищенню рівня політичної культури та громадянської активності студентів.
        \end{enumerate}

    1.4.2. ЦВКс взаємодіє з органами студентського самоврядування (Студентською радою Київського авіаційного інституту (СР КАІ), Студентською радою студмістечка (СР СМ), Студентською уповноваженою делегацією (СУД), студентськими радами факультетів/інститутів (СРФ/СРІ)), адміністрацією Університету та студентськими організаціями у порядку, визначеному Положенням про ОСС та цим Положенням.

\subsection*{1.5. Термін повноважень}
\addcontentsline{toc}{subsection}{1.5. Термін повноважень}
    1.5.1. Строк повноважень складу ЦВКс становить 1 (один) рік з дня обрання КСУ.

    1.5.2. Повноваження ЦВКс можуть бути достроково припинені за рішенням КСУ у випадках, передбачених цим Положенням.

    1.5.3. У разі дострокового припинення повноважень ЦВКс, КСУ протягом одного місяця обирає новий склад комісії на строк до завершення звичайного терміну повноважень.